\section{Introduction and Context}
\label{sec:introduction}

This section establishes the conceptual foundation for understanding SI-Mapper.
It explains what building ontologies are, why a standardized vocabulary for building systems matters,
and why the ASHRAE 223P standard was chosen as the target output format.
Readers do not need a technical background to follow this discussion.

\subsection{The Semantic Metadata Gap in Buildings}
\label{subsec:metadata-gap}

Every commercial building contains hundreds of sensors, actuators, and control points.
A modern office block may have thousands of temperature sensors, pressure transducers, damper position
indicators, and setpoint registers --- each one a source of valuable operational data.
Yet despite this wealth of instrumentation, the data is largely inaccessible to software that was
not purpose-built for the specific building it operates in.

The reason is a fundamental naming problem.
Each Building Management System (BMS) vendor uses its own proprietary conventions for labelling control
points.
A temperature sensor named \texttt{2500.AI11} in one system becomes \texttt{RTU-3/SA-T} in another.
Neither label carries enough information for software to determine automatically what physical equipment
the point belongs to, what role the measurement plays, or how that equipment connects to other components.
Without that context, building software cannot be deployed without a building-specific customization step.

As the United States Department of Energy states directly, ``building software is not `plug and play' ---
technicians must manually map sensors and actuators to software inputs through tedious point mapping
processes that vary building-by-building.''
This manual, building-by-building effort is the central barrier this report addresses.

The solution is a \emph{standardized metadata layer} --- a machine-readable description of what every
sensor and actuator means, which equipment it belongs to, and how that equipment connects to the rest of
the system.
This is what building ontologies provide.

\subsection{Building Ontologies: A Standardized Language for Buildings}
\label{subsec:ontologies}

A building ontology is a shared, machine-readable vocabulary that describes building systems, their
components, and how they connect.
Think of it as a dictionary that any software application can consult: if a sensor is declared to be a
supply air temperature measurement on the outlet of a fan coil unit, any compliant application
can understand that without manual configuration.

Four major building ontologies have emerged over the past decade.
Each reflects different design priorities and is backed by a different community.

\begin{table}[htbp]
\centering
\caption{The four major building ontology standards.}
\label{tab:ontologies}
\small
\begin{tabular}{lllll}
\toprule
\textbf{Ontology} & \textbf{Governed By} & \textbf{Approach} & \textbf{Focus} & \textbf{Maturity} \\
\midrule
Project Haystack & Haystack Connect & Tag-based dictionary & Equipment types, sensor tags & Widely deployed, informal \\
Brick Schema & Open consortium & RDF/OWL semantic graph & Equipment, sensors, spatial hierarchy & Production ready (v1.4) \\
RealEstateCore & RealEstateCore Consortium & OWL ontology & Property management + technical systems & Production ready \\
ASHRAE 223P & ASHRAE + DOE/NREL & Formal OWL ontology & Topology, connections, telemetry & Second advisory public review \\
\bottomrule
\end{tabular}
\end{table}

\textbf{Project Haystack} is the oldest and most widely deployed approach.
It uses a tag-based system where sensors are described by a set of keywords (e.g., \texttt{air temp sensor}).
Its informality makes it easy to adopt but difficult to enforce consistently across projects.

\textbf{Brick Schema} builds formal semantic rules on top of a Haystack-compatible vocabulary.
It uses the RDF/OWL knowledge representation standard, making it more consistent and machine-interpretable
than Haystack.
Brick is production-ready and in active deployment across research institutions and building operators.

\textbf{RealEstateCore} focuses on the property management perspective --- occupancy, asset management,
and energy performance --- with significant overlap into technical building systems.
It is used mainly in the European commercial real estate sector.

\textbf{ASHRAE 223P} is the most rigorous of the four.
It defines not just equipment types but full system topology: how equipment connects, what flows between
components, and how those connections link to live sensor data.
ASHRAE's BACnet committee, Project Haystack, and Brick Schema are actively collaborating to integrate
their vocabularies into ASHRAE 223P as the unifying standard.

A 2025 systematic comparison of these four ontologies found that no single standard covers all needs
in practice.
A complete building model is expected to be ``a mix of ASHRAE 223P, Brick, QUDT, RealEstateCore, and
possibly additional standards.''
ASHRAE 223P is designed as the semantic backbone that ties the others together.

\subsection{ASHRAE Standard 223P}
\label{subsec:ashrae-223p}

ASHRAE Standard 223P is a proposed ASHRAE standard that formally defines a knowledge representation
framework for building systems.
Its goal is to enable software applications to determine the meaning and context of building data
without manual, building-specific configuration.

\textbf{Current status:} The standard is in its second advisory public review as of 2025--2026.
It is not yet ratified.
The project receives \$1.3 million in Department of Energy funding and is led by the National Renewable
Energy Laboratory (NREL), with participation from Lawrence Berkeley National Lab, Pacific Northwest
National Lab, NIST, Colorado School of Mines, and Cornell University.

ASHRAE 223P organizes building system information around six core concepts, described in
Table~\ref{tab:223p-concepts}.

\begin{table}[htbp]
\centering
\caption{The six core concepts of ASHRAE 223P.}
\label{tab:223p-concepts}
\small
\begin{tabular}{lp{5.5cm}p{5cm}}
\toprule
\textbf{Concept} & \textbf{Definition} & \textbf{Example} \\
\midrule
Type & Well-defined classes with formal rules governing usage & Fan, Pump, CoolingCoil \\
Topology & How equipment and spaces connect & Fan outlet $\to$ Duct $\to$ Coil inlet \\
Composition & Hierarchical grouping of entities & AHU contains Fan, Coil, Damper \\
Telemetry & Observable properties with external data references & Supply air temp $\to$ BACnet object 2500.AI11 \\
Characteristics & Descriptive attributes with units & Rated airflow in m\textsuperscript{3}/s \\
Medium & Substance conveyed through connections & Chilled water, Supply air, Electricity \\
\bottomrule
\end{tabular}
\end{table}

\textbf{Topology} is the most distinctive concept in ASHRAE 223P.
Where other ontologies describe what equipment exists, 223P also describes how equipment connects.
Every connection has a direction (inlet, outlet, or bidirectional), a physical path (duct, pipe, or
wire), and a medium (air, water, electricity).
This means an ASHRAE 223P model captures not just ``there is a fan and a coil'' but ``air flows from
the fan outlet through a duct segment into the coil inlet, conveying Supply Air at design conditions.''

\textbf{Telemetry and external references} allow live BACnet sensor data to be explicitly linked to
semantic properties in the model.
A supply air temperature sensor in the model can declare that its value comes from
\texttt{BACnet object 2500.AI11} --- closing the loop between the abstract ontology and the live
building management system.

\textbf{Why SI-Mapper chose ASHRAE 223P:} It is the only standard that simultaneously captures full
HVAC topology, links BACnet points to semantic equipment properties, and is explicitly designed for
advanced building controls and grid services.
Its backing by DOE, NREL, and ASHRAE's standardization process makes it the most credible long-term
choice for interoperability.

\subsection{Advantages and Barriers of Semantic Building Models}
\label{subsec:advantages-barriers}

Creating a semantic building model is not a trivial investment.
Understanding why organizations would make that investment --- and why they currently do not ---
frames the business case for automation tools like SI-Mapper.

\textbf{Advantages of semantic building models:}

\begin{itemize}[leftmargin=*, itemsep=2pt]
  \item \textbf{Software portability:} applications built on a standard ontology can work across
        buildings without building-specific customization, just as web browsers work across websites.
  \item \textbf{Advanced analytics:} fault detection, automated commissioning, and predictive
        maintenance all require knowing not just sensor values but what those values mean and how
        equipment relates to each other.
  \item \textbf{Grid services enablement:} Virtual Power Plant (VPP) and Non-Wire Alternative (NWA)
        programs require machine-readable building load data; semantic models provide this in a
        standardized form.
  \item \textbf{Reduced integration costs:} once a semantic model exists, any compliant application
        can consume it without a custom integration project.
  \item \textbf{Digital twin support:} the semantic layer is the connective tissue between physical
        building systems and their digital representations.
\end{itemize}

\textbf{Barriers to adoption:}

\begin{itemize}[leftmargin=*, itemsep=2pt]
  \item \textbf{High creation cost:} building a semantic model today requires specialized domain
        expertise and significant labor --- often weeks of work by a qualified HVAC engineer.
  \item \textbf{No automated tooling for existing buildings:} most available approaches are manual
        or semi-manual, requiring expert knowledge of both HVAC systems and ontology standards.
  \item \textbf{Fragmented standards landscape:} organizations must choose between partially
        overlapping standards with no single dominant solution.
  \item \textbf{Model maintenance:} models drift from reality as building systems are modified,
        requiring ongoing expert effort to keep them current.
  \item \textbf{Low industry adoption:} despite clear benefits, the broader industry has been slow
        to adopt formal ontologies, creating a chicken-and-egg problem for software vendors.
\end{itemize}

The high cost of \emph{creating} semantic models --- not deploying them --- is the central barrier.
A semantic model that would take a qualified engineer several weeks to create manually could enable
years of automated analytics and grid integration.
The investment is worthwhile; the bottleneck is the creation step.
This is the problem SI-Mapper is designed to solve.
